\section{Le Chat: inferencia rápida, capa de herramientas y volante de retroalimentación del producto}

En la entrevista, el producto de consumo cumple tres funciones: mostrar las capacidades del modelo y de las herramientas, generar una distribución real de tareas y recolectar retroalimentación sobre los fallos. La clase también señala que una inference optimization extrema puede reducir la frecuencia de actualización del modelo o la compatibilidad de arquitectura. Es un trade-off de sistemas típico: la velocidad no es gratuita, pues puede exigir un runtime fijo, un kernel especializado, una familia de modelos restringida o un proceso de publicación más complejo.

\subsection{Por qué el producto forma parte del sistema de entrenamiento}

Las solicitudes de los usuarios, los tool trace y las calificaciones explícitas pueden revelar problemas que el benchmark no ve: qué consultas necesitan web search, qué tareas de código fallan con frecuencia y en qué idiomas o sectores el desempeño es inestable. Después, el equipo de producto convierte esos fallos en evaluation slice y en prioridades de entrenamiento, con lo que se forma un data flywheel.

\lecturefigure{12-product-feedback-flywheel.png}{Los productos del tipo de Le Chat conectan el serving, las herramientas y las prioridades de entrenamiento en un volante de retroalimentación.}{Subtítulos locales 24:45--27:40; redibujo conceptual.}

\begin{knowledgebox}{Lectura de la figura: la retroalimentación primero debe convertirse en una failure taxonomy}
Un simple thumbs up/down tiene un significado ambiguo: una calificación baja puede deberse a un error factual, al tono, a que una herramienta agotó su tiempo de espera, a la latencia de la interfaz o a un cambio en el objetivo del usuario. Un volante confiable debe conservar el request context, las versiones del modelo y de las herramientas, la latencia, el resultado final y la clasificación manual, y solo entonces decidir si se modifican los datos, el modelo, el prompt, la tool o el serving.
\end{knowledgebox}

\begin{warningbox}{«Recolectar datos del producto» requiere límites de gobernanza}
Los registros del producto pueden contener información personal, secretos empresariales o datos regulados. Reincorporarlos al entrenamiento requiere purpose limitation, retention policy, control de acceso, elección del usuario, desidentificación y lineage de los datos. El volante de retroalimentación no consiste en agregar automáticamente todas las conversaciones al conjunto de entrenamiento.
\end{warningbox}

\teachervoice{La clase también presenta una señal concreta del producto: en cierto momento, las tareas de código representaban una proporción muy alta de las solicitudes en inglés, por lo que el equipo aumentó la inversión en code generation. Este ejemplo muestra que el roadmap no proviene solo del criterio de los directivos, sino que también puede derivarse de la distribución observable de las tareas; aun así, los datos observados deben corregirse por el sesgo de la población de usuarios.}

\subsection{Resumen de la sección}

El producto es a la vez una interfaz de servicio, un muestreador de tareas y un puesto de observación de fallos. Un volante de retroalimentación de alta calidad debe descomponer las señales de los usuarios en problemas localizables y devolverlas a la evaluación y al entrenamiento bajo una gobernanza de la privacidad.
